\chapter{Analisi dei requisiti}\label{ch:analysis}

In questo capitolo verranno descritti i requisiti funzionali e non funzionali del sistema.
\section{Requisiti funzionali}
A seguire i requisiti funzionali individuati dalla richiesta:
\begin{itemize}
    \item aggiunta e rimozione dei clienti: l'utente deve poter aggiungere e rimuovere clienti nel sistema, rappresentati come nodi del grafo;
    \item aggiornamento costo spostamento tra clienti: l'utente deve poter aggiornare il peso degli archi, che rappresenta il costo di spostarsi da un cliente all'altro.
            Questo costo può tenere conto sia della condizione dei percorsi, che il costo in termini di tempo e denaro;
    \item visualizzazione algoritmo ACO in tempo reale: l'utente deve poter visualizzare le approssimazioni del cammino migliore, mentre l'algoritmo è in esecuzione;
    \item uso di grafi completi: è stato richiesto che l'implementazione usi un grafo completo, anche se quello fornito
            dal cliente non lo sia. In questo caso il grafo viene completato con archi \textit{fittizi} con pesi sufficientemente elevati. Non si vuole che i pesi siano infiniti ma neppure che siano così bassi che possano venire presi al posto di un arco con peso \textit{reale}.
\end{itemize}

\subsection{Descrizione delle funzioni}
\begin{itemize}
    \item Aggiunta dei clienti
        \begin{itemize}
            \item Input: tripla formata da identificatore del cliente, identificatore della sede relativa a questo, lista dei costi per spostarsi dal cliente ad altri clienti.
            \item Output: conferma da parte del sistema dell'aggiunta, o messaggio di errore nel caso in cui l'operazione non sia andata a buon termine.
            \item Effetto: aggiunta del nodo cliente all'interno della memoria locale del nodo ACO che fa riferimento alla sede \texttt{Sede}, con relativi archi.
        \end{itemize}
    \item Rimozione dei clienti
        \begin{itemize}
            \item Input: identificatore del cliente da rimuovere.
            \item Output: conferma da parte del sistema della rimozione, o messaggio di errore nel caso in cui l'operazione non sia andata a buon termine.
            \item Effetto: rimozione del nodo cliente dalla memoria locale del nodo ACO che lo conteneva e rimozione archi che comprendono il cliente.
        \end{itemize}
    \item Aggiornamento costo spostamento tra clienti
        \begin{itemize}
            \item Input: lista di triple formate da i due clienti e il costo dello spostamento tra loro.
            \item Output: conferma della da parte del sistema dell'aggiornamento, o messaggio di errore nel caso in cui l'operazione non sia andata a buon termine.
            \item Effetto: aggiornamento dei costi degli archi per ogni coppia \texttt{(Cliente1, Cliente2)} nei relativi nodi ACO.
        \end{itemize}
    \item Visualizzazione grafo
        \begin{itemize}
            \item Input: nessuno
            \item Output: grafo aggiornato in tempo reale
            \item Effetto: visualizzazione da parte dell'utente del grafo aggiornato in tempo reale.
        \end{itemize}
\end{itemize}

\section{Requisiti non funzionali}

È stato deciso di utilizzare il linguaggio Erlang per l'implementazione del sistema.

Le trasparenze del sistema sono come descritte nella sezione~\ref{sec:tra}.

Non sono stati richiesti limiti per quanto concerne prestazioni o tempistiche del tempo, ma è stato chiesto
di usare una versione asincrona dell'algoritmo ACO, se possibile. Non sono state effettuate richiesta rispetto alla sicurezza del sistema.
A tale scopo, si è individuata una variante asincrona, presentata in~\cite{ILIE2013762}, come punto di partenza.

I clienti sono identificati da ID unici e non è possibile aggiungere multipli clienti con lo stesso ID.